\documentclass[11pt]{article}
\usepackage[margin=1in]{geometry}
\usepackage[T1]{fontenc}
\usepackage[utf8]{inputenc}
\usepackage{amsmath,amssymb,amsthm}
\usepackage{enumitem}

\title{ICPC World Finals 2025\\D. Buggy Rover}
\author{}
\date{}

\begin{document}
\maketitle

\section*{Problem Summary}

The rover always follows one of the $4!=24$ possible direction orderings. For each recorded move we
know the current cell and the direction actually taken, but the ordering may change between consecutive
moves. We must minimize how many times the ordering changed.

\section*{Key Observations}

\begin{itemize}[leftmargin=*]
    \item There are only $24$ possible direction orderings, so we can treat each ordering as a small DP state.
    \item For a fixed move step and a fixed ordering, it is easy to test whether that ordering is compatible
    with the recorded move: the chosen direction must be valid, and every valid direction that appears
    earlier in the ordering must be absent.
    \item Once we know, for every step, which of the $24$ orderings are allowed, the rest is a shortest-path
    problem on a layered graph:
    staying with the same ordering costs $0$, switching to a different ordering costs $1$.
\end{itemize}

\section*{Algorithm}

\begin{enumerate}[leftmargin=*]
    \item Enumerate all $24$ permutations of $\{N,E,S,W\}$.
    \item Simulate the recorded walk once to recover the rover position before each move.
    \item For every step and every permutation, test whether that permutation could have produced the
    recorded move from that position, and store the result in a boolean table \texttt{allowed[step][perm]}.
    \item Run dynamic programming over the move sequence:
    \[
    dp[i][p] = \text{minimum number of changes after processing moves }0..i
    \]
    with permutation $p$ used at move $i$.
    The transition is
    \[
    dp[i][p] = \min_q \bigl(dp[i-1][q] + [p \ne q]\bigr),
    \]
    restricted to allowed states.
\end{enumerate}

\section*{Correctness Proof}

We prove that the algorithm returns the correct answer.

\paragraph{Lemma 1.}
For any step and any permutation, \texttt{allowed[step][perm]} is true exactly when that permutation
could have produced the recorded move at that step.

\paragraph{Proof.}
The rover chooses the first direction in its ordering that leads to a valid neighboring cell. Therefore a
permutation is compatible exactly when the recorded direction is valid and every other valid direction that
appears earlier in the permutation is impossible. This is exactly the test performed by the program. \qed

\paragraph{Lemma 2.}
The DP value $dp[i][p]$ equals the minimum possible number of ordering changes among all explanations
of the first $i+1$ moves that use permutation $p$ on move $i$.

\paragraph{Proof.}
The base case $i=0$ is correct because there is no earlier move, so every allowed first permutation costs
$0$ changes. For $i>0$, any valid explanation ending with permutation $p$ on move $i$ must come from
some permutation $q$ used on move $i-1$. If $q=p$ we pay no extra cost; otherwise we pay exactly one
change. Taking the minimum over all valid $q$ gives precisely the recurrence above. \qed

\paragraph{Theorem.}
The algorithm outputs the correct answer for every valid input.

\paragraph{Proof.}
By Lemma 1, the DP only considers permutations that are locally consistent with the recorded moves.
By Lemma 2, each DP value is the minimum number of changes among all globally consistent explanations
with the specified final permutation. Therefore the minimum over the last layer is exactly the minimum
number of times the rover's direction ordering could have changed. \qed

\section*{Complexity Analysis}

Let $m=|s|$. Testing all step-permutation pairs costs $O(24 \cdot 4 \cdot m)=O(m)$. The DP costs
$O(24^2 m)$, which is well within the limits. The memory usage is $O(24m)$ for the compatibility table
and $O(24)$ for the rolling DP arrays.

\section*{Implementation Notes}

\begin{itemize}[leftmargin=*]
    \item The rover position must be updated using the recorded move sequence, not by replaying guessed
    permutations.
    \item A two-row rolling DP is enough because the transition only depends on the previous step.
\end{itemize}

\end{document}
